\chapter*{Resumen}
\addcontentsline{toc}{chapter}{Resumen} % Añade "Resumen" a la Tabla de Contenido

Generar el pulso que una aplicación requiere de un láser ultrarrápido equivale a fijar la
configuración de parámetros de su cavidad, un problema de optimización de caja negra, no
convexo y multiestable frente al cual los métodos de ajuste tradicionales resultan
insuficientes. La literatura ha respondido con algoritmos inspirados en la naturaleza y
técnicas de aprendizaje automático, pero deja dos vacíos: se ha concentrado casi por
completo en el régimen de \emph{mode-locking} pasivo, dejando el activo subexplorado, y
está organizada como una colección de estudios de caso aislados, sin comparaciones que
mantengan fijos el láser, el espacio de parámetros y la función objetivo para variar solo
el algoritmo. Esta tesis aborda ambos vacíos con un estudio comparativo controlado de tres
estrategias libres de gradiente: un algoritmo genético (GA), un híbrido con red neuronal
sustituta (NN+GA) y la optimización por enjambre de partículas (PSO), sobre un simulador de
la ecuación maestra de Haus validado contra soluciones analíticas, en los regímenes activo
y pasivo, bajo una misma función de mérito y métricas comunes de calidad, costo y robustez.

Sobre una matriz de sesenta corridas independientes y contraste estadístico no paramétrico,
las tres estrategias alcanzan soluciones factibles, pero el método preferible depende del
régimen. En el régimen activo el PSO domina de forma significativa en
calidad, dispersión, velocidad de convergencia y robustez del óptimo, asociada a un
acortamiento solitónico más fuerte; en el pasivo el PSO y
el GA quedan estadísticamente empatados. El sustituto neuronal iguala al GA en el régimen
activo con $5{.}7$ veces menos evaluaciones del simulador y sacrifica calidad en el pasivo,
aunque un control emparejado muestra que, con ese presupuesto, su aporte no se distingue
del de su propio muestreo inicial. Un único factor (la forma
del paisaje de optimización de cada régimen, regular en el activo y multiestable en el
pasivo) explica los tres hallazgos, y la robustez local emerge como una dimensión de
comparación que discrimina entre métodos cuando la calidad los iguala. El trabajo, realizado
en simulación, aporta un marco reproducible y evidencia de que el régimen activo es, para
esta familia de métodos, el más tratable; la validación experimental y la extensión al
aprendizaje por refuerzo profundo quedan como líneas de trabajo futuro.

\vspace{1em}
\noindent\textbf{Palabras clave:} láseres ultrarrápidos, \emph{mode-locking}, optimización
sin gradiente, algoritmo genético, optimización por enjambre de partículas, redes neuronales
sustitutas, ecuación maestra de Haus, aprendizaje automático.
